\textbf{Sea-state sensing.}
The canonical instrument for sea-state monitoring is the moored wave buoy, which records heave acceleration and integrates it to recover sea-surface elevation. Buoys yield accurate spectral and directional estimates but require fixed deployment infrastructure and return only a single spatial point. X-band marine radar~\cite{rs9121261} is the established ship-mounted alternative: backscatter images of the sea surface encode the directional wave spectrum and can be processed in near-real-time, but the hardware is heavy, power-hungry, and designed for vessels far larger than the WIG platforms targeted here. Single-point laser and ultrasonic altimeters~\cite{toa2020ultrasonic} are lightweight and compact but similarly provide no spatial coverage. The WIG application requires a sensing solution that is lightweight, operates from a moving low-altitude platform, and reconstructs the sea surface geometrically ahead of the vehicle -- a combination that none of these established modalities addresses directly.

\textbf{Stereo reconstruction of the sea surface.}
Non-contact stereo wave measurement was first demonstrated at oceanographic scale by Benetazzo~\cite{benetazzo2006offshore} and subsequently productised in the open-source WASS pipeline~\cite{bergamasco2017wass}, which achieved sub-centimetre reconstruction accuracy from a near-fixed, near-vertical platform using Hirschmüller's semi-global matching (SGM)~\cite{hirschmuller2008stereo}. Later work targeting the sea surface's photometric adversity -- low texture over large areas and specular non-Lambertian reflections that violate brightness constancy -- includes GPU-accelerated feature-based stereo on ship-mounted rigs~\cite{app12157447}, stereo for maritime object localisation~\cite{jmse12010197}, and thermal stereo~\cite{li2025thermal}, which sidesteps visible-spectrum specularity by imaging in the long-wave infrared at the cost of a more specialised and expensive sensor pair.

\textbf{Stereo-matching architectures.}
Classical SGM~\cite{hirschmuller2008stereo} applies global path-wise smoothness regularisation over the disparity image, allowing it to recover coherent depth over mildly textureless regions -- making it the strongest classical baseline available without training data. Learned architectures have consistently outperformed it on standard benchmarks: HITNet~\cite{DBLP:journals/corr/abs-2007-12140} propagates disparity hypotheses over a hierarchical tile pyramid without building a full 4D cost volume, enabling real-time inference; RAFT-Stereo~\cite{DBLP:journals/corr/abs-2109-07547} applies iterative GRU updates over a multi-level correlation volume for high accuracy at higher latency; WAFT-Stereo~\cite{wang2026waftstereowarpingalonefieldtransforms} replaces cost-volume lookups with warping-alone field transforms, reaching state-of-the-art benchmark performance at latency comparable to RAFT-Stereo.

\textbf{Wave-parameter estimation and platform motion.}
Spectral wave parameters -- significant wave height $H_s = 4\sqrt{m_0}$, peak period $T_p$, and spectral mean periods $T_{m01}$/$T_{m02}$ -- are defined against the directional wave spectrum, whose canonical empirical forms are the Pierson--Moskowitz spectrum for fully-developed seas~\cite{pierson1964spectral} and the JONSWAP spectrum for fetch-limited wind-sea growth~\cite{Hasselmann1973}. Altimetric estimation from a moving platform introduces the encounter-frequency effect: a platform travelling through the wave field at speed $U$ measures an encounter period rather than the true period, the two being related via the deep-water dispersion relation and the relative heading between platform and waves~\cite{dean1991water}. Ichikawa et al.~\cite{ichikawa2024seasurface} demonstrated this correction explicitly for a UAV-mounted nadir LiDAR altimeter -- the closest published configuration to the one used in this project -- and quantified the resulting period shift.

\textbf{Short-term wave forecasting.}
Deterministic wave forecasting from reconstructed wave fields propagates the phase-resolved spectrum forward in time using linear wave theory~\cite{naaijen2010real,qi2018predictable}, with the forecast horizon set by the predictable zone of the measurement. Data-driven methods using recurrent networks~\cite{hochreiter1997long} and attention-based architectures~\cite{vaswani2017attention} have been applied to longer-horizon statistical forecasting. This work focuses on parameter \emph{estimation} -- a prerequisite for any forecasting stage -- and treats forecasting as a descoped but natural next step (Section~\ref{sec:future}).
